
\subsection{The Ideal $\leftrightarrow$ Noisy Replacement Argument}

A notation that we will need later is the extending of the infinity characterization to slices, 
\begin{definition}[Inifinite Slice]
  \label{def:infislice}
  Let $K$ be a slice in a Toom's contour. We will say that $K$ is infinite slice if there is $g\in S$ and vertices $u,v$ such: 
  \begin{equation*}
    \begin{split}
      L_{g}(v) - L_{g}(u) \ge n 
    \end{split}
  \end{equation*}
\end{definition}


\inb{\textbf{Consider to change it to a claim. Propagated error has not introduced yet.}}
\begin{remark}
Note that $u$ and $v$ might be set on different connected chains. In particular, the origin of $K$ vertices might be a finite (yet disconnected) chain, and yet $K$ might be an infinite slice. In addition, if $K$ is a slice in the last layer of the analysis Toom's contour at step $\tau$, and $P_{t^{\star} -1}$ is the reduced propagated error up to that step, then we have a transition in the opposite direction, namely if $K$ is finite then the propagated error $P_{t^{\star}}$ is finite.
\end{remark}


  \inb{\textbf{Change.}} 
  Our construction will result (after initialization) in a Toric encoded circuit, where all the quantum registers are encoded using the $(d,d^{\prime})$-Toric code induced by a Toric with a side of length $n$. Then, for a decoder $D$, we are going to take the following simple procedure: just apply the sweep-cycle $  n^{d+1}  $ times - the number of vertices times the side's length. If we want him to correct a measurement, namely a classical state which is the outcome of the measurement, then we perform only the first stage of the sweep-cycle (without moving into the dual code). 

The following claim states that for ideal running (not subjected to noise), assuming that during the decoding all the slices are finite, the decoding succeeds in correcting them.
\begin{claim}
  \label{claim:decodestopnoise}
Let $\ket{\psi}$ be a logical codeword of the $(d,d^{\prime})$-Toric code with side length $n$. Let $P \in \mathcal{P}^{(\cdot)}$. Assuming that during the application of $n^{d+1}$ ideal sweep-cycles over the $P\ket{\psi}$, the Toom's counter does not admit an infinite slice, then at the end the error is erased, namely the result of the computation is $\ket{\psi}$.
\end{claim}
\begin{proof}
We will prove that the decoder erases the $X$-component of $P$, and the proof for the $Z$-component is identical. Let $v_{1}, v_{2}, \ldots, v_{m}$ be the vertices supporting the faces of $X(P)$. During the ideal computation, no new faults are introduced, and therefore $v_{1}, \ldots, v_{m}$ are the only leaves in the induced Toom's contour.
  Moreover, they set a natural relation over its slices as follows: two slices $K_{1}$ and $K_{2}$, not necessarily with the same time coordinate, will be said to be \textit{history-related}, $K_{1}\sim K_{2}$, if there is a leaf $v_{i}$ that belongs to both of their histories. Thus, for a time $t$ and a leaf $v \in v_{1}, \ldots, v_{m}$, we can define the notation $K_{v,t}$ to be the slice at time $t$ that is reached from leaf $v$ by arrows. We emphasis that $\sim$ is not an equivalence relation, since it doesn't satisfy transitivity. 
Yet, let us match slices to classes when the map is many-to-many. The delegate of the slices reached from vertex $v$, namely the slices $\{ K_{v,t} : t \}$, will be denoted by $\mathbf{K}_{v}$. Proving that at the end of time $n^{d+1}$, for any leaf $v$, the slice $K_{v,t}$ is empty, or equivalently is of size zero, implies the correctness of the claim.


  For proving this, let us present the following terminology. For leaves $v,u,w$, and a time coordinate $t$, consider the following: 
  \begin{enumerate}
    \item Suppose that $K_{v,t}$ and $K_{u,t+1}$ are slices at times $t$ and $t+1$, such $K_{v,t}\sim K_{u,t+1}$. If $\Size{K_{u,t+1}} > \Size{K_{v,t}}$, we will say that the class $\mathbf{K}_{v}$ grows at time $t$. Similarly, if $\Size{K_{u,t+1}} < \Size{K_{v,t}}$, then we will say that slice $\mathbf{K}_{v}$ shrank at time $t$.
    \item Suppose that $K_{v,t}$ and $K_{u,t}$ are slices at time $t$, and $K_{w,t+1}$ is a slice at time $t+1$, such that $K_{v,t}\not\sim K_{u,t}$, $K_{w,t+1}\sim K_{v,t}$, and $K_{w,t+1}\sim K_{u,t}$. Then we will say that $\textbf{K}_{v}$ and $\textbf{K}_{u}$ merged at time $t$.
    \item The size of a class $\mathbf{K}_{v}$ at time $t$ is $\Size{K_{v,t}}$.
  \end{enumerate}
Since no new faults are introduced during the computation, we have that in each step, the only option for a class to grow is by merging. Notice that, by definition, classes cannot 'split'. Namely, if $K_{v,t}\sim \mathbf{K}_{u}$, then for any time $t^{\prime} \ge t$, it holds that $K_{v,t^{\prime}} \sim \textbf{K}_{u}$. Specifically, a correction by the sweep rule that disconnects two error clusters that were connected by checks and, prior to correction, belonged to the same slice and therefore to the same class, would still belong to the same class after the correction.


Therefore, there can be at most $m-1$ merges. And since the number of classes is bounded by the number of vertices on the Toric, we have that $m \le n^{d}$, namely at most $n^{d}$ merges can occur. Hence, during the $n^{d+1}$ sweep-cycles, there has to be a continuous sequence of $n$ cycles for which no merge occurs. Since it's given that the size of all the slices is lower than $n$, the sizes of all connected errors are also lower than $n$. So, a sweep-cycle decreases the size of a finite error by at least 1 (\Cref{claim:polesmove}). Thus, at the end of that sequence of $n$ cycles, all the slices' sizes shrink to zero; namely, the connected errors are erased after those cycles.
\end{proof}


\inb{\textbf{Explantation has not introduced yet.}}

To prove that no slice reaches a size greater than $n$, we use explanation trees. In particular, an explanation tree that characterizes the running well is one that, prior to decoding, allows the introduction of new faults, and therefore leaves, and during the ideal $n^{d+1}$ iterations, no new leaf is introduced. The way to analyze the probability of a proper explanation occurring is to bound it by the probability\footnote{Actually, by an upper bound of the probability} of the same explanation occurring in the setting where the following $n^{d+1}$ iterations are also subjected to local stochastic noise.
 
\begin{claim}[Replacement argument]
  \label{claim:replace}
  Consider a running of Toric encoded circuit with depth $l$. Denote by $A$ the event in which, after an additional $ n^{d+1}$ \textbf{ideal} (non-noisy) iterations of the sweep-cycle, the error is not erased, and denote by $B$ the event that, during the application of an additional $n^{d+1}$ \textbf{noisy} iterations, the Toom's contour admits an infinite slice. Let $f:(n,d,d^{\prime},l,p) \rightarrow \mathbb{R} $ be an upper bound on $\prb{B}$ that assumes only that the noise is distributed according to a local stochastic noise channel with noise rate $p$. Then, we have:
  \begin{equation*}
    \begin{split}
      \prb{A} \le f(n,d,d^{\prime},l,p)
    \end{split}
  \end{equation*}
\end{claim}
\begin{proof} 
\Cref{claim:decodestopnoise} states that if during the $ n^{d+1}$ ideal sweep-cycles the Toom's contour does not admit an infinite slice, then the error is erased. Consider the set of atomic events corresponding to the faults in the first $l$ noisy layers, such that during the additional $ n^{d+1}$ ideal iterations there is a step that admits an infinite slice. All those atomic events are also contained in $B$, since local stochastic noise is our only assumption. It's possible that noise in the $B$ scenario is distributed such that the probability of a fault occurring during the additional $ n^{d+1}$ iterations is zero. In that specific case, $\prb{A} = \prb{B}$ and therefore $\prb{A} \le f(n,d,d^{\prime},l,p)$.
\end{proof}

